\section{Discussion and Limitations}
\label{sec:discussion}

\subsection{Statistical Uncertainty of DES Time-Averaging}
\label{sec:des_statistics}

A fundamental assumption in the pipeline is that each DES run produces a well-converged time-mean $\overline{C_d}$ that can be used as a training observation for the correction GP. This assumption is examined by computing the effective number of independent samples in the averaging window.

Adjacent force coefficient samples at $\Delta t = 4\times10^{-4}$~s are strongly correlated: the integral time scale $\tau_{C_d}$ estimated from the autocorrelation function of the time series is $\approx12$~ms for $C_d$ and $\approx10$~ms for $C_l$. Over an averaging window of $T_\mathrm{avg} = 0.32$~s, this yields:

\begin{equation}
    N_\mathrm{eff} = \frac{T_\mathrm{avg}}{2\,\tau} \approx \frac{0.32}{0.024} \approx 13 \quad (C_d)
    \label{eq:neff}
\end{equation}

independent observations, rather than the $813$ raw samples. The resulting 95\% confidence interval on the time-mean $\overline{C_d}$ per case is:

\begin{equation}
    \sigma_{\overline{C_d}} = \frac{1.96\,\sigma_{C_d}}{\sqrt{N_\mathrm{eff}}} \approx \frac{1.96\times0.013}{\sqrt{13}} \approx \pm0.007
    \label{eq:cdci}
\end{equation}

For $C_l$ the situation is substantially worse: $\sigma_{\overline{C_l}} \approx \pm0.035$ per case, comparable to or larger than the inter-case $C_l$ differences the surrogate is attempting to learn.

Three practical consequences follow. First, the minimum detectable $C_d$ difference between two cases at 95\% confidence is $2\times0.007 = 0.014$; the surrogate cannot reliably distinguish designs that differ by less than this. Second, surrogate-predicted $C_l$ differences below $\sim\!0.07$ should not be interpreted as significant. Third, the correction GP's WhiteKernel noise term, fitted during hyperparameter optimisation, should absorb the $\pm0.007$ per-case noise if the kernel is expressive enough---but this requires $N_\mathrm{eff}$ to be representative of the true noise distribution, which is not guaranteed when the DES cases are run at different parameter settings with different separated flow topologies.

The recommended fix is to extend $t_\mathrm{end}$ from $0.45$~s to $1.0$~s for future campaigns. This quadruples the averaging window to $0.87$~s ($\approx33$ effective samples per case), halving $\sigma_{\overline{C_d}}$ to $\approx\pm0.004$ and making individual cases sufficient to detect differences of $\sim\!0.008$.

\subsection{EI Non-Convergence and the Flat Optimum}
\label{sec:flat_optimum}

The EI plateau at $\approx$0.005--0.01 across 20 BO iterations indicates a structural limitation of BO in this setting, not a failure of the acquisition function. The Expected Improvement scales as $\mathrm{EI} \propto \sigma(\mathbf{x}_*)$ near the optimum: if the GP has a long length scale (broad uncertainty) and a shallow objective landscape, each new observation reduces $\sigma$ locally but reveals comparable uncertainty nearby, maintaining a roughly constant EI.

In physical terms, the high-slant regime ($\alpha_\mathrm{slant} > 35^\circ$) produces a slowly varying vortex topology where the aerodynamic response to geometry changes is broad and smooth rather than sharply peaked. The BO therefore spread evaluations across this region without converging to a single point, which is correct behaviour — the objective landscape does not have a sharp global minimum to converge onto.

From a practical standpoint, EI non-convergence is not fatal: the GP surrogate correctly identified the global trend and confirmed case~021 ($\alpha_\mathrm{slant}=32.5^\circ$) as the incumbent through all 20 BO iterations. The MF co-Kriging correction (\S\ref{sec:mf_results}) then relocates the optimum accounting for the L2--L3 fidelity gap.

\subsection{Multi-Fidelity Correction and Optimum Verification}
\label{sec:verification_gap}

The initial single-fidelity campaign identified $\alpha_\mathrm{slant}=32.5^\circ$, $\alpha_\mathrm{diff}=8.9^\circ$ as the L2 RANS optimum ($f = 0.210$). As quantified in \S\ref{sec:l3_validation}, a single L3 validation run at this point revealed a fidelity gap of $\Delta C_l = +0.160$, inflating the L3 objective to $f = 0.248$ — an 18\% degradation. The root cause is the coarser L2 boundary layer delaying slant separation onset and artificially strengthening the trailing vortex system.

To close this gap, an additive co-Kriging correction was fitted to 8 L3 anchor simulations spanning the full design space (\S\ref{sec:l3_anchors}). The corrections $\delta_{C_l}$ at the anchor locations range from $-0.114$ (at $40.0^\circ/12.0^\circ$) to $+0.167$ (at $21.3^\circ/15.4^\circ$), with a mean of $+0.042$ and standard deviation $0.087$. The corrected HF surrogate relocates the objective minimum accounting for this spatially varying offset, and the resulting MF optimum is verified with a dedicated L3 solve (Table~\ref{tab:mf_optimum}, \S\ref{sec:mf_results}).

The key methodological insight is that the fidelity gap is not a constant offset: $\delta_{C_l}$ varies spatially with $\alpha_\mathrm{slant}$ because the degree of slant separation — and hence the L2 mesh's over-prediction of vortex strength — changes across the design space. Near the critical angle ($21.3^\circ$--$32.5^\circ$) the correction exceeds $+0.12$; at high slant ($38.9^\circ$) where the flow is fully separated the correction is only $+0.012$. A constant-bias correction (applying the single observed $\Delta C_l = +0.160$ at the old optimum everywhere) would mis-locate the corrected minimum by over-correcting in the high-slant basin where the MF optimum actually resides. The co-Kriging model learns this spatial variation from the 8 anchors and propagates it correctly to the re-optimisation.

With 8 anchor points in a 2-D space the correction GP is deliberately conservative: the WhiteKernel noise term absorbs anchor-to-anchor variation that may reflect physical flow-regime changes rather than GP model error. The predictive uncertainty $\sigma_\mathrm{HF}$ is therefore wider than the L2 surrogate alone. The verification L3 solve at the MF optimum ($\alpha_\mathrm{slant}=38.9^\circ$, $\alpha_\mathrm{diff}=10.5^\circ$) returned $f_1 = 0.1975$, against the surrogate prediction of $f_1 = 0.1964$ — a $0.6\%$ error, well within the predicted $\pm 2\sigma_\mathrm{HF}$ band. This confirms both the surrogate's fidelity at the optimum and the validity of the $+4.9^\circ$ diffuser shift as a genuine feature of the HF landscape rather than a surrogate artefact.

\subsection{RANS Turbulence Model Limitations}
\label{sec:rans_limits}

The $k$-$\omega$~SST model used for the RANS DoE is known to under-predict mixing in separated shear layers, producing the positive $\delta_{C_d}$ correction values observed in Table~\ref{tab:l3_anchors_results}. More critically, the RANS solution is steady-state: near the critical slant angle ($\alpha_\mathrm{slant} \approx 28$--$32^\circ$), the actual flow periodically switches between attached and separated states, and the steady RANS converges to an averaged intermediate state that does not correspond to either physical regime. This means RANS DoE data in the critical region is structurally mis-specified, not merely noisy, and the co-Kriging correction GP must learn a large, non-smooth $\delta$ in that zone.

The consequence for surrogate quality is that the correction GP's length scale in the slant direction is shorter near the critical angle, correctly increasing surrogate uncertainty there and widening the $\sigma_\mathrm{HF}$ band. The MF optimum's uncertainty is therefore physically motivated: the correction GP is uncertain precisely in the region where the RANS model is most unreliable.

\subsection{Simplified Geometry Relative to Full Race Car}
\label{sec:ahmed_limits}

The Ahmed body is a canonical benchmark, not a race car. Three simplifications introduce limitations when interpreting results in a motorsport context:

\textbf{No moving ground plane.} The simulations use a stationary ground boundary condition, consistent with the Lienhart \& Becker \cite{lienhart2002} wind-tunnel setup. Real ground-effect flows require a moving ground (tangential velocity matching $U_\infty$), which strengthens the underbody venturi effect and changes the separation behaviour at the diffuser exit. The stationary ground underestimates diffuser efficiency.

\textbf{No rotating wheels.} Wheel rotation generates significant additional drag and modifies the ground plane boundary layer upstream of the diffuser. The cylindrical leg geometry used here is a simplified support structure, not a rotating tyre.

\textbf{No cooling flows.} Real ground vehicles pass significant mass flow through radiators, which exits through body outlets and interacts with the underbody and rear wake. The Ahmed body is solid.

These simplifications mean that the absolute coefficient values ($C_d \approx 0.27$--$0.29$ at the optimum) are not directly transferable to a race car. The optimisation trends and the MF framework itself generalise fully to more complex geometries.
